% from ega4-2-fr.tex
\oldpage[IV-2]{145}

\subsection{Propriétés de transitivité}

\begin{proposition}[6.6.1]
\label{IV.2.s6.6.1-fr}
Pour un préschéma localement noethérien $Z$, désignons par $P(Z)$ l'une
quelconque des propriétés suivantes:
\begin{enumerate}
\item[(a)] $Z$ est un préschéma de Cohen--Macaulay.
\item[(b)] $Z$ vérifie $(S_n)$.
\item[(c)] $Z$ est régulier.
\item[(d)] $Z$ vérifie $(R_n)$.
\item[(e)] $Z$ est normal.
\item[(f)] $Z$ est réduit.
\end{enumerate}
Soient alors $X$, $Y$, $Z$ trois préschémas localement noethériens,
$f:X\to Y$, $g:Y\to Z$ deux morphismes, et posons $h=g\circ f$.
\begin{enumerate}
\item[(i)] Supposons que $f$ et $g$ soient plats et que, pour tout
$y\in f(X)$ (resp. tout $z\in g(Y)$), le préschéma $f^{-1}(y)$
(resp. $g^{-1}(z)$) vérifie $P$. Alors $h$ est plat et, pour tout
$z\in h(X)$, $h^{-1}(z)$ vérifie $P$.
\item[(ii)] Supposons vérifiées les conditions suivantes: $f$ est
fidèlement plat, $h=g\circ f$ est plat, pour tout $y\in Y$, le
préschéma $f^{-1}(y)$ vérifie $P$, et, pour tout $z\in h(X)$, le
préschéma $h^{-1}(z)$ vérifie $P$. Alors $g$ est plat et, pour tout
$z\in g(Y)$, $g^{-1}(z)$ vérifie $P$.
\end{enumerate}
\begin{enumerate}
\item[(i)] On sait déjà (2.1.6) que $h$ est plat; d'autre part, pour
      tout $z\in Z$,
\[
 f_z=f\otimes 1_{k(z)}:h^{-1}(z)\longrightarrow g^{-1}(z)
\]
est plat (2.1.4), et, pour tout $y\in g^{-1}(z)$, $f_z^{-1}(y)$ est
isomorphe à $f^{-1}(y)$ par transitivité des fibres (I, 3.6.4). La
conclusion résulte alors respectivement de (6.3.5, (ii)), (6.4.2),
(6.5.2, (ii)), (6.5.3, (ii)), (6.5.4, (ii)) et (3.3.5, (ii)).
\item[(ii)] On sait déjà que $g$ est plat (2.2.13); en outre, pour
      tout $z\in Z$, $f_z$ est fidèlement plat (2.2.13). La conclusion
      résulte alors respectivement de (6.3.5, (i)), (6.4.2),
(6.5.2, (i)), (6.5.3, (i)), (6.5.4, (ii)) et (2.1.13).
\end{enumerate}
\end{proposition}

\begin{remark}[6.6.2]
\label{IV.2.s6.6.2-fr}
Supposons $h$ plat, $f$ fidèlement plat, et supposons que, pour tout
$z\in Z$, $h^{-1}(z)$ soit de coprofondeur $\leqslant n$ (5.7.1); alors
il résulte du raisonnement de (6.6.1, (ii)) et de (6.3.2) que
$g^{-1}(z)$ est de coprofondeur $\leqslant n$ pour tout $z\in g(Y)$ et
que, pour tout $y\in Y$, $f^{-1}(y)$ est de coprofondeur $\leqslant n$.
\end{remark}

